\section{Attribution and verification}
The two constructions, their symmetry assertions and their use in the finite
normal-pair theorem are due to D\'iaz Ramos, Molinier and Viruel. The
restriction-and-extension criterion is elementary classical reasoning; its
explicit separation from recognition and rigidity is an organization of the
proofs, not a priority claim. Lemma~\ref{UR.Components} isolates the component
argument also used in the recursive application. No external graph
representation theorem is needed to make the proofs above complete.

The accompanying \texttt{RecognizableRestriction.lean} constructs the
restriction homomorphism from its actual action on an embedded vertex set,
proves its kernel description, derives injectivity from pointwise rigidity and
surjectivity from existence of extensions, constructs the group equivalence,
and transports a distinguished subgroup. It does not formalize the incidence
or degree arguments proving those hypotheses for any gadget. The complete
chamber, cleaning and component arguments above remain written proofs.

The two complete GAP programs in the separate appendix provide bounded
checks, including exceptional graphs and small normal pairs. Their computations
do not replace any quantified theorem above. GAP and GRAPE's exact finite
automorphism algorithms remain part of the trust boundary of those checks;
Lean uses its kernel and the pinned mathlib infrastructure. Neither the number
of programs nor the number of formal declarations measures mathematical
compression.

\begin{thebibliography}{9}
\bibitem{DMV} A. D\'iaz Ramos, R. Molinier and A. Viruel,
\emph{Every finite group is represented by a finite incidence geometry},
arXiv:2511.02410v2, 16 December 2025.
\end{thebibliography}
